\documentclass[nonacm,sigconf]{acmart}
\usepackage{paralist}
\usepackage{xurl}
\settopmatter{printacmref=false}
% TODO: Replace the descriptive artifact title and anonymous author with
% the submission's exact title and author metadata, as appropriate.
\title{EdgeTPU Firmware Recovery and Fault-Oracle Analysis Artifacts}
%\author{Anonymous Authors}
\renewcommand{\shortauthors}{Anonymous Authors}
% TODO: Supply the archive DOI, license, paper section/table references,
% measured person/compute times, and reference recovery/scanner results.
\newcommand{\artifactfile}[1]{\path{#1}}
\begin{document}
\maketitle
\appendix
\section{Artifact Appendix}

\subsection{Abstract}
The artifact provides a three-stage EdgeTPU firmware-analysis workflow:
acquiring memory dumps on a prepared Pixel~9a, importing dumps into IDA
and recovering candidate functions, and scanning reconstructed firmware
for candidate fault-oracle memory-access gadgets. The package includes
acquisition source and binaries, two reference logs, a firmware image,
IDAPython scripts, and a prepared IDA Pro~9.4 database. Reviewers can
start from either the device, supplied logs, or prepared database.
The package's redistribution license remains to be specified; IDA and
Hex-Rays are proprietary dependencies and are not bundled.

\subsection{Description \& Requirements}
The root \artifactfile{README.md} documents all commands and expected
output formats. \artifactfile{POC/DumpAuto_stable/} contains
\artifactfile{runner}, its source, the \artifactfile{cracktpu} acquisition
helper, and \artifactfile{compile.sh}. \artifactfile{Public_artifacts/}
contains the firmware, dump logs, importer, function-recovery script,
and scanner. The prepared database is
\artifactfile{edgetpu-rio_eval_re.fw.i64} at the package root.

\subsubsection{Security, privacy, and ethical concerns}
Acquisition requires root, stops the EdgeTPU logging service, and can
reboot the device or interrupt SSH. Use the prepared review device or
an authorized test device. After reconnecting, re-enter root and stop
logging before resuming; restore logging afterward. Import and recovery
modify the database, so use a working copy. Scanner reports are
candidates, not automatically confirmed vulnerabilities.

\subsubsection{How to access}
Use the submitted package. For remote acquisition, follow the Tailscale
invitation in README Section~1.1 and connect your computer to the reviewer
network. Log in to the prepared Pixel~9a using
\texttt{ssh -p 8022 u0\_a292@100.78.159.78}
and enter \texttt{septlab} when prompted for the SSH password.
Root access is already enabled: after logging in, run \texttt{su} to
enter a root shell; no additional rooting steps are required.
Contact the authors through the artifact-review channel if the invitation
expires or the device is offline. Offline stages need no device connection.

\subsubsection{Hardware dependencies}
Live acquisition uses the authors' prepared, rooted Pixel~9a with access
to \artifactfile{/dev/edgetpu-soc}. Other devices or firmware builds
are not assumed compatible. Offline evaluation uses a workstation
capable of running IDA Pro~9.4; minimum RAM and disk requirements have
not been measured.

\subsubsection{Software dependencies}
IDA Pro~9.4 with IDAPython is required; Stage~3 also needs a licensed
ARM Hex-Rays decompiler. Remote access requires Tailscale and SSH.
Rebuilding and local deployment require Android NDK and ADB.
The build script uses \texttt{darwin-x86\_64} and
\artifactfile{aarch64-linux-android34-clang}; adjust the NDK host path on
other operating systems. SSH access does not imply ADB access.

\subsubsection{Benchmarks}
Inputs are \artifactfile{Public_artifacts/edgetpu-rio.fw}, the prepared
database, and two logs:
\artifactfile{0x0ff00000_0x0ff0E4FF_dump.log} (14,656 words; 2,175 tagged)
and \artifactfile{0x10000000_0x1000ea80_dump.log} (15,009 words; 2,150 tagged).
Records contain an address and a 32-bit value, optionally marked
\texttt{cache\_unstable}. These counts describe bundled inputs; unstable
locations need not match byte for byte in a fresh acquisition.

\subsection{Set Up}
\subsubsection{Installation}
Extract the package and install the dependencies above. For device
acquisition, use the deployed binaries or set \texttt{NDK} and run
\texttt{sh compile.sh} in \artifactfile{POC/DumpAuto_stable/} with an
ADB-connected device. After logging in to the prepared Pixel via SSH,
run \texttt{su} to enter a root shell, then run
\texttt{stop edgetpu\_logging}; enter \artifactfile{/data/local/tmp/},
make both helpers executable, and create \artifactfile{review-dumps/}.
For a fresh IDA import, load the raw firmware as ARM little-endian at
\textbf{base address \texttt{0x0ffff000}}. Alternatively, open a copy
of the prepared \texttt{.i64} directly; do not rebase or reimport it.

\subsubsection{Basic test}
For the shortest offline check, open the prepared database, wait for
analysis, and execute \artifactfile{Public_artifacts/faultoracle_scanner_mc.py}
via \texttt{File > Script file}. The IDA output window should show
per-handler reports and end with
\texttt{TOTAL sites (after per-def-site dedup): ...}.
Check for \texttt{no function at} or \texttt{gen\_microcode failed};
a final total alone does not prove that every handler was analyzed.
No verified reference site count is supplied in the README.

\subsection{Evaluation Workflow}
\subsubsection{Major claims}
The following workflow claims follow the README's result-to-artifact
mapping. Exact paper section/table references and quantitative acceptance
criteria must be supplied before claiming reproduction of the paper's
main results.
\begin{compactitem}
\item[(C1)] The acquisition workflow recovers the two memory regions
and supports resumption after interruption (E1).
\item[(C2)] Imported dump words support firmware reconstruction and
heuristic recovery of candidate function boundaries (E2).
\item[(C3)] The microcode scanner identifies candidate fault-oracle
memory-access sites for manual inspection (E3).
\end{compactitem}

\newpage
\subsubsection{Experiments}
Measured human effort, runtime, and disk usage for E1--E3
are not provided in the README and remain to be supplied. E1 includes
repeated reads and may require reconnection; no runtime bound is claimed.

\begin{compactitem}
\item[(E1)] \textbf{Memory acquisition (C1; README Stage~1).}
\emph{Preparation:} Connect, enter root, and stop logging as above.
\emph{Execution:} In \artifactfile{/data/local/tmp/}, invoke
\texttt{./runner} \texttt{START} \texttt{END} \texttt{./cracktpu} \texttt{OUTPUT} twice: first with
\artifactfile{START=0x0ff00000}, \artifactfile{END=0x0ff0e500}; then with
\artifactfile{START=0x10000000}, \artifactfile{END=0x1000ea84}. Use the matching log
names under \artifactfile{review-dumps/}. Endpoints are exclusive and
steps are four bytes. After reboot, reconnect, repeat \texttt{su} and
\texttt{stop edgetpu\_logging}, and rerun the same command/output path.
Runner resumes after the last valid record. Finish with
\texttt{start edgetpu\_logging}.
\emph{Results:} Retrieve logs and check their final word addresses:
\texttt{0x0ff0e4fc} and \texttt{0x1000ea80}, respectively. Inspect
instability tags and incomplete acquisitions rather than requiring
identical values at every location.

\item[(E2)] \textbf{Import and recovery (C2; README Stage~2).}
\emph{Preparation:} Load the firmware at the stated base and save a copy.
\emph{Execution:} Run \artifactfile{import_dump_log_to_ida.py} twice,
selecting one supplied log each time; wait for reanalysis. Run
\artifactfile{recover_tpu_funcs.py}, checking that \texttt{SEG\_NAME}
(default \texttt{ROM}) names the intended segment. Repeat for other
relevant segments, such as \texttt{.dump\_0ff}, as needed.
\emph{Results:} Expect 14,656 and 15,009 patched words if all addresses
are mapped; investigate malformed/unmapped records. Save recovery output
ending in \texttt{DONE --- added ... funcs total} and inspect candidate
functions. Zero additions may reflect prior recovery or a wrong segment.

\item[(E3)] \textbf{Fault-oracle scanning (C3; README Stage~3).}
\emph{Preparation:} Use the prepared database and ARM decompiler.
\emph{Execution:} Run the scanner as in the basic test; it analyzes
configured \texttt{HANDLERS} at \texttt{MMAT\_CALLS} maturity.
\emph{Results:} Save IDA/Hex-Rays versions, database identity, and
per-handler output, including site addresses and deduplicated counts.
Inspect reported instructions and pseudocode. Sites are fault-oracle
candidates, not conventional ROP gadgets or proof of exploitability.
\end{compactitem}

\subsection{Notes on Reusability}
Bundled logs and the prepared database allow offline evaluation without
the Pixel. New firmware versions require reviewing segment names,
handler addresses, and payload registers. Function recovery is heuristic;
scanner failures must not be interpreted as absence of candidates.
A permanent archive, license, measured costs, and paper-specific reference
results remain necessary release/evaluation metadata.

\subsection{Version}
Based on the LaTeX template for artifact evaluation version 1.0.

%% AUTHORSHIP AND PROVENANCE OF THIS TEMPLATE

%% This template is based on CCS 2025 template. The AEC Chair for CCS 2026 was:
%% Jérémie Decouchant, Delft University of Technology
%% Fengwei Zhang, Southern University of Science and Technology
%% Jun Dai, Worcester Polytechnic Institute

%% This template was created by David Balenson and Ning Zhang for the ACM CCS 2025
%%
%% The AEC Chair for CCS 2025 was:
%% David Balenson (USC Information Science Institute)
%% Ning Zhang (Washington University in St. Louis)

%% This template is based on CCS 2024 template, 
%%  https://www.sigsac.org/ccs/CCS2024/assets/appendix/ccs2024-ae.zip
%%  the AEC chair for CCS 2024 was 
%%  Eric Eide, University of Utah
%% the template was originally based on the "USENIX Security 2023 AE Template."
%% https://secartifacts.github.io/usenixsec2023/appendix/usesec23-ae-latex.zip
%%
%% The AEC Co-chairs for USENIX Security 2023 were:
%% Cristiano Giuffrida (Vrije Universiteit Amsterdam)
%% Anjo Vahldiek-Oberwagner (Intel Labs)

%% The USENIX Security AE template was based on the
%% % Artifact Appendix Template for EuroSys'22 AE
%% https://sysartifacts.github.io/eurosys2022/appendix/EuroSys22_ArtifactAppendix_template.tex
%%
%% The AEC Co-chairs for EuroSys 2022 were:
%% Natacha Crooks (UC Berkeley)
%% Solal Pirelli (EPFL)
%% Salvatore Signorello (University of Lisbon)
%% Anjo Vahldiek-Oberwagner (Intel Labs)

%% The above were based on:

%% LaTeX template for Artifact Evaluation V20220119
%%
%% Original Authors
%% * Grigori Fursin (cTuning foundation, France) 2014-2020
%% * Bruce Childers (University of Pittsburgh, USA) 2014
%%
%% Modified by
%% * Clementine Maurice (CNRS, France) 2021-2022
%% * Cristiano Giuffrida (Vrije Universiteit Amsterdam, Netherlands) 2021-2022
%%
%% (C) Copyright 2014-2022
%%
%% CC BY 4.0 license
%%




\end{document}
\endinput
